% Procedencia: RESULTADO_RECUPERADO del capítulo 24 de la integral de 2249
% páginas y de FORMA_ELIPTICA_RADIO_MODULAR_ACCION_ORIENTADA.md (D04).
% Delta expositivo: dominio del operador diagonal escrito completamente y
% recuperación del entero en las dos cartas cocientadas hecha explícita.
% Dependencia externa de etiqueta: sec:ellipse. Bibliografía: elipse,
% hmtedicionintegral, flm, borcherds. No requiere macros nuevas.

\section{De l’ellipse d’action à la dualité T}
\label{sec:ii-dualidad-t}

La réciprocité de l’ellipse fournit une coordonnée de forme pour
l’échange entre valeur visible et mémoire de frontière. L’antécédent est le
même état enrichi d’APP--TRIT--TPK : APP conserve les feuillets de
somme et de produit avec leurs résidus et quotients ; TRIT conserve leur orientation ;
TPK transporte le feuillet, la retenue, le chemin et la mémoire. Les sorties angulaires et
d’action déjà construites sur la structure discrète conjointe du continu
déterminent ensuite l’ellipse. Cette section compose ses lecteurs ultérieurs :
forme elliptique, échange réticulaire, réalisation de quantité de mouvement--enroulement
et évaluation modulaire \cite{elipse,hmtedicionintegral}.

Le point de coupure se situe donc après la génération de
\(A,C^*,\hbar\). Aucun rayon prescrit ni aucune constante de corde ne sélectionne
ces sorties. La mémoire de frontière est d’abord conservée dans la paire
résidu--quotient, puis transportée au moyen des applications démontrées
ci-après.

\subsection{La forme d’action détermine le rayon normalisé}

Reprenons la coordinatisation de la section~\ref{sec:ellipse}. Soient \(A\ne0\),
\(\hbar>0\) et \(|C^*/A|<1\), les deux angles étant exprimés dans la même
unité. Les demi-axes restent étiquetés :
\begin{equation}
 \eta_{\rm el}=\operatorname{artanh}(C^*/A),\qquad
 a_S=\sqrt{2\hbar}\,e^{\eta_{\rm el}/2},\qquad
 b_S=\sqrt{2\hbar}\,e^{-\eta_{\rm el}/2}.
 \label{eq:ii-t-elipse}
\end{equation}
Les coordonnées équilibrées de phase et leurs demi-axes ont pour unité la racine
d’action. Le rayon qui suit est adimensionnel.

\begin{proposition}[Rayon de forme et réciprocité étiquetée]
\label{prop:ii-t-radio}
La coordinatisation rectangulaire de l’ellipse détermine de manière unique
\begin{equation}
 R_0=\sqrt{\frac{b_S}{a_S}}
 =e^{-\eta_{\rm el}/2}
 =\left(\frac{A-C^*}{A+C^*}\right)^{1/4}>0,
 \qquad \tau_{\rm el}=iR_0^2.
 \label{eq:ii-t-radio}
\end{equation}
Elle conserve le rapport angulaire signé au moyen de
\begin{equation}
 \frac{C^*}{A}=\frac{1-R_0^4}{1+R_0^4}.
 \label{eq:ii-t-recuperacion-angular}
\end{equation}
L’échange des demi-axes étiquetés transforme
\((\eta_{\rm el},R_0,\tau_{\rm el})\) en
\((-\eta_{\rm el},R_0^{-1},-1/\tau_{\rm el})\), en laissant fixe
l’échelle \(a_Sb_S=2\hbar\).
\end{proposition}

\begin{proof}
La division des demi-axes de~\eqref{eq:ii-t-elipse} donne
\(b_S/a_S=e^{-\eta_{\rm el}}\). L’application \(R\mapsto R^2\)
est bijective sur \(\mathbb R_{>0}\), de sorte que sa racine positive est
unique. Pour \(x=C^*/A\in(-1,1)\),
\(2\operatorname{artanh}x=\log((1+x)/(1-x))\) ; donc
\(R_0^4=(1-x)/(1+x)\). Isoler \(x\) prouve
\eqref{eq:ii-t-recuperacion-angular}. Échanger \(a_S,b_S\)
inverse leur rapport et change le signe de \(\eta_{\rm el}\). Enfin,
\(-1/(iR_0^2)=iR_0^{-2}\), tandis que le produit des demi-axes
reste égal à \(2\hbar\).
\end{proof}

Le rapport angulaire est retrouvé ; les deux angles absolus et leur histoire
génératrice continuent de résider dans l’état enrichi. Le point \(R_0=1\)
correspond à la forme circulaire \(C^*/A=0\), non à l’identification de
toutes les orientations ou de toutes les histoires.

\subsection{Échange réticulaire et équivalence d’opérateurs}

\begin{definition}[Domaine stable de la lecture duale]
\label{def:ii-t-reticular}
Sur \(\mathscr D_0=\mathbb Z^2\times\mathbb R_{>0}\), nous définissons
\begin{equation}
 E_{R_0}(m,w)=\frac{m^2}{R_0^2}+w^2R_0^2,\qquad
 \mathsf T_0(m,w,R_0)=(w,m,R_0^{-1}).
 \label{eq:ii-t-energia}
\end{equation}
La première coordonnée prolonge la lecture visible et la seconde celle du
quotient. Le domaine entier complet permet de les échanger sans exiger
que le quotient appartienne à une section de neuf représentants.
\end{definition}

Pour le rayon de~\eqref{eq:ii-t-radio}, la matrice de cette forme est
\begin{equation}
 D(\eta_{\rm el})=
 \begin{pmatrix}e^{\eta_{\rm el}}&0\\0&e^{-\eta_{\rm el}}\end{pmatrix}
 =\frac{1}{2\hbar}
 \begin{pmatrix}a_S^2&0\\0&b_S^2\end{pmatrix}.
 \label{eq:ii-t-matriz}
\end{equation}
En effet, \(R_0^{-2}=e^{\eta_{\rm el}}\) et
\(R_0^2=e^{-\eta_{\rm el}}\). L’équation de l’ellipse
\(Q^2/a_S^2+P^2/b_S^2=1\), multipliée par \(2\hbar\), utilise
la matrice inverse \(D(-\eta_{\rm el})\). Ainsi sont différenciées
la matrice des demi-axes au carré et celle de l’équation elliptique.

\begin{theorem}[Dualité réticulaire et conjugaison unitaire]
\label{thm:ii-t-unitaria}
La transformation \(\mathsf T_0\) est involutive et satisfait
\begin{equation}
 E_{R_0^{-1}}(w,m)=E_{R_0}(m,w).
 \label{eq:ii-t-invariancia}
\end{equation}
Dans \(\mathcal H_{\rm lat}=\ell^2(\mathbb Z^2)\), soit \(H(R_0)\)
l’opérateur diagonal de valeurs \(E_{R_0}(m,w)\), de domaine
\[
 \mathcal D(H(R_0))=
 \left\{c\in\ell^2(\mathbb Z^2):
 \sum_{m,w\in\mathbb Z}E_{R_0}(m,w)^2|c_{m,w}|^2<\infty\right\}.
\]
Alors \(H(R_0)\) est autoadjoint et non négatif. L’échange
\(U_T|m,w\rangle=|w,m\rangle\) est une involution unitaire qui
transporte les domaines et satisfait
\begin{equation}
 U_T H(R_0)U_T^{-1}=H(R_0^{-1}).
 \label{eq:ii-t-conjugacion}
\end{equation}
\end{theorem}

\begin{proof}
Échanger deux fois les entiers et inverser deux fois le rayon ramène
au point initial. De plus,
\[
 E_{R_0^{-1}}(w,m)=w^2R_0^2+m^2R_0^{-2}=E_{R_0}(m,w).
\]
Le domaine diagonal contient les suites à support fini et est dense.
La multiplication par une suite réelle est symétrique sur ce domaine.
Si \(c\) appartient au domaine de l’adjoint, l’évaluation contre chaque
vecteur de la base impose que l’adjoint ait pour coordonnées
\(E_{R_0}(m,w)c_{m,w}\). Son appartenance à \(\ell^2\) équivaut
exactement à la condition qui définit \(\mathcal D(H(R_0))\).
L’opérateur et son adjoint coïncident donc. Ses valeurs diagonales non
négatives prouvent la positivité.

L’échange permute bijectivement la base orthonormée, il est donc
unitaire et son inverse est lui-même. Pour toute suite \(c\),
\[
 \sum_{m,w}E_{R_0^{-1}}(m,w)^2|(U_Tc)_{m,w}|^2
 =\sum_{m,w}E_{R_0}(m,w)^2|c_{m,w}|^2.
\]
Cette égalité prouve le transport exact des domaines. Enfin,
\((U_TH(R_0)U_T^{-1}c)_{m,w}=E_{R_0}(w,m)c_{m,w}
=E_{R_0^{-1}}(m,w)c_{m,w}\) ; ainsi est démontrée
\eqref{eq:ii-t-conjugacion} sur tout le domaine, non seulement sur la base.
\end{proof}

L’orientation ajoute une information que l’énergie ne distingue pas. Dans le plan
de phase, l’échange \(S_2(Q,P)=(P,Q)\) et la rotation
\(J_2(Q,P)=(-P,Q)\) transportent la même forme quadratique, mais
\(S_2^*(dQ\wedge dP)=-dQ\wedge dP\) et
\(J_2^*(dQ\wedge dP)=dQ\wedge dP\). Avec
\(\lambda=P\,dQ\), on obtient directement
\[
 S_2^*\lambda=d(PQ)-\lambda,\qquad
 J_2^*\lambda=\lambda-d(PQ).
\]
L’intégration sur une courbe fermée donne des actions opposées sous \(S_2\)
et égales sous \(J_2\). Par conséquent, l’échange des indices de
\eqref{eq:ii-t-energia} et le transport symplectique d’une trajectoire
sont des applications distinctes, même lorsqu’ils publient la même énergie.

\subsection{Sections nonadiques et transport de la retenue}

Pour \(n\in\mathbb Z\), la section positive écrit
\[
 r_9^+(n)=1+((n-1)\bmod9),\qquad
 Q_9^+(n)=\frac{n-r_9^+(n)}9,\qquad n=r_9^+(n)+9Q_9^+(n).
\]
Le reste de \(n-1\) est pris dans \(\{0,\ldots,8\}\). Le passage à
la section des restes \(0,\ldots,8\) est
\begin{equation}
 C(r,Q)=\bigl(r-9\varepsilon,Q+\varepsilon\bigr),\qquad
 \varepsilon=\begin{cases}1,&r=9,\\0,&r\ne9.\end{cases}
 \label{eq:ii-t-acarreo}
\end{equation}
Ainsi, \(r+9Q\) reste fixe. Pour \(n=9\), les deux coordonnées
sont \((9,0)\) et \((0,1)\) : leurs énergies non corrigées sont
\(81/R_0^2\) et \(R_0^2\). L’égalité de l’entier n’implique pas
l’égalité de ces insertions énergétiques. La forme covariante suivante
résout le changement de section sans supprimer la retenue.

\begin{definition}[Deux coordinatisations quotientées de la lecture nonadique]
\label{def:ii-t-cociente}
Soient \(L_9=\mathbb Z(9,-1)\), \(S(m,w)=(w,m)\) et
\(L_9^\vee=SL_9=\mathbb Z(-1,9)\). Pour
\(L\in\{L_9,L_9^\vee\}\), nous définissons
\[
 \mathcal E_{R_0}^{L}([x]_L)
 =\min_{\ell\in L}E_{R_0}(x+\ell).
\]
Le transport entre coordinatisations est
\[
 \mathsf T_{\rm cov}:
 (\mathbb Z^2/L_9)\times\mathbb R_{>0}
 \longrightarrow
 (\mathbb Z^2/L_9^\vee)\times\mathbb R_{>0},\qquad
 ([x]_{L_9},R_0)\longmapsto([Sx]_{L_9^\vee},R_0^{-1}).
\]
\end{definition}

\begin{theorem}[Covariance avec mémoire de franchissement]
\label{thm:ii-t-cociente}
Les énergies et le transport précédents sont bien définis. Le changement
de section~\eqref{eq:ii-t-acarreo} conserve la classe dans
\(\mathbb Z^2/L_9\), et
\begin{equation}
 \mathcal E_{R_0^{-1}}^{L_9^\vee}([Sx]_{L_9^\vee})
 =\mathcal E_{R_0}^{L_9}([x]_{L_9}).
 \label{eq:ii-t-covariancia}
\end{equation}
L’échange en retour retrouve la classe et le rayon initiaux. Les
applications de coordonnées entières \(N([x])=x_1+9x_2\) et
\(N^\vee([y])=9y_1+y_2\) sont des isomorphismes avec \(\mathbb Z\)
et satisfont \(N^\vee([Sx])=N([x])\).
\end{theorem}

\begin{proof}
Remplacer \(x\) par \(x+\ell_0\), avec \(\ell_0\in L\),
permute les candidats au minimum. Celui-ci existe : pour \(L=L_9\),
\[
 E_{R_0}(x+k(9,-1))
 =\left(\frac{81}{R_0^2}+R_0^2\right)k^2
  +2\left(\frac{9x_1}{R_0^2}-x_2R_0^2\right)k+E_{R_0}(x),
\]
ce qui tend vers l’infini avec \(|k|\) ; pour \(L_9^\vee\), le
coefficient dominant est \(R_0^{-2}+81R_0^2>0\). Le changement
\(C(x)-x=-\varepsilon(9,-1)\) appartient à \(L_9\).
Si \(x'-x\in L_9\), alors \(Sx'-Sx\in L_9^\vee\),
ce qui prouve que le transport est bien défini. Par
\eqref{eq:ii-t-invariancia},
\[
 \min_{\ell\in L_9}E_{R_0^{-1}}(Sx+S\ell)
 =\min_{\ell\in L_9}E_{R_0}(x+\ell).
\]
Cela démontre~\eqref{eq:ii-t-covariancia} ; \(S^2=I\) prouve
le retour entre les deux coordinatisations.

Enfin, l’homomorphisme \(x\mapsto x_1+9x_2\) est
surjectif et son noyau est exactement \(L_9\), car
\(x_1=-9x_2\). De même, le noyau de
\(y\mapsto9y_1+y_2\) est \(L_9^\vee\). Leurs applications
induites sur les quotients sont donc des isomorphismes. Substituer
\(y=Sx\) prouve l’identité finale.
\end{proof}

L’entier retrouvé détermine ses résidus et quotients une fois la
section fixée. Le représentant qui minimise l’énergie peut ne pas être ce
représentant canonique. Le quotient est une représentation covariante de la
lecture nonadique ; l’état source conserve en outre le feuillet, le chemin, l’orientation,
la frontière, la mémoire et la survie. On ne remplace pas non plus \(L_9^\vee\)
par \(L_9\) : l’échange a lieu entre deux coordinatisations transportées.

\subsection{Quantité de mouvement, enroulement et échelle de la réalisation}

La réalisation circulaire introduit une coordinatisation dimensionnelle explicite.
Notons \(R>0\) son rayon et \(\alpha'>0\) un paramètre
de dimension longueur au carré. Ce symbole est distinct de la sortie
adimensionnelle \(\alpha_{\rm HMT}\). Fixer la coordinatisation n’altère ni le
générateur angulaire ni le rayon de forme. La normalisation et son inverse sont
\[
 \mathcal N_{\alpha'}(m,w,R)
 =\left(m,w,\frac{R}{\sqrt{\alpha'}}\right),\qquad
 R=\sqrt{\alpha'}\,R_0.
\]
Les indices se lisent maintenant comme nombre de quantité de mouvement et nombre
d’enroulement ; les composantes du mode zéro, exprimées en unités
d’inverse de longueur, sont
\begin{equation}
 p_L=\frac mR+\frac{wR}{\alpha'},\qquad
 p_R=\frac mR-\frac{wR}{\alpha'},\qquad
 E_{\rm str}=\frac{p_L^2+p_R^2}{2}
 =\frac{m^2}{R^2}+\frac{w^2R^2}{\alpha'^2}.
 \label{eq:ii-t-quirales}
\end{equation}

\begin{theorem}[Conjugaison dimensionnelle de la dualité T]
\label{thm:ii-t-cuerda}
La transformation
\[
 \mathsf T_{\alpha'}(m,w,R)=\left(w,m,\frac{\alpha'}R\right)
\]
est une involution et satisfait
\begin{equation}
 \mathcal N_{\alpha'}\mathsf T_{\alpha'}
 =\mathsf T_0\mathcal N_{\alpha'},\qquad
 E_{\rm str}=\alpha'^{-1}E_0\circ\mathcal N_{\alpha'},
 \label{eq:ii-t-normalizacion}
\end{equation}
où \(E_0(m,w,R_0)=E_{R_0}(m,w)\). Elle conserve \(E_{\rm str}\),
laisse \(p_L\) fixe et inverse le signe de \(p_R\). Pour des données
oscillatoires \(N_L,N_R\) et des intercepts \(a_L,a_R\) fixes,
elle conserve également la condition
\begin{equation}
 \frac{\alpha'}4(p_L^2-p_R^2)+N_L-N_R-a_L+a_R=0,
 \quad\text{c’est-à-dire}\quad
 N_L-N_R=a_L-a_R-mw.
 \label{eq:ii-t-nivel}
\end{equation}
\end{theorem}

\begin{proof}
Deux applications échangent deux fois les indices et envoient
\(R\) sur \(\alpha'/(\alpha'/R)=R\). De même,
\[
 \mathcal N_{\alpha'}\mathsf T_{\alpha'}(m,w,R)
 =\left(w,m,\frac{\sqrt{\alpha'}}R\right)
 =\mathsf T_0\mathcal N_{\alpha'}(m,w,R).
\]
La substitution dans la forme normalisée donne
\[
 E_0\left(m,w,\frac R{\sqrt{\alpha'}}\right)
 =\frac{\alpha'm^2}{R^2}+\frac{w^2R^2}{\alpha'}
 =\alpha'E_{\rm str}(m,w,R).
\]
Pour \(R'=\alpha'/R\), les modes transformés sont
\[
 p'_L=\frac{wR}{\alpha'}+\frac mR=p_L,\qquad
 p'_R=\frac{wR}{\alpha'}-\frac mR=-p_R.
\]
Par conséquent, tant la somme de leurs carrés que leur différence sont conservées.
Le développement direct donne \(p_L^2-p_R^2=4mw/\alpha'\), ce qui
transforme la première équation de~\eqref{eq:ii-t-nivel} en la seconde.
L’échange conserve \(mw\) ; les autres données étant fixes, il conserve
la condition complète.
\end{proof}

La formule abrégée \(N_L-N_R=nw\) exige, avec ces orientations,
\(a_L=a_R\) et \(n=-m\). Maintenir le signe de
\eqref{eq:ii-t-nivel} conserve le dictionnaire entre les deux lectures.
Le résultat démontré comprend la conjugaison des rayons, les modes
zéro, l’énergie et la compatibilité des niveaux déclarée. L’interprétation
dimensionnelle identifie les coordonnées au moyen de leur coordinatisation ; elle ne transforme pas
\(R_0\) à lui seul en longueur et ne déduit pas une théorie des champs
complète d’une égalité d’énergies.

\subsection{L’évaluation modulaire et le caractère de Moonshine}

Le corridor exceptionnel déjà développé conserve l’incidence du même
antécédent jusqu’à \(V^\natural\). Son caractère gradué a pour expression
\(J(\tau)=j(\tau)-744\), avec l’invariance modulaire employée dans
Moonshine \cite{flm,borcherds}. Cette \(J\) est une fonction sur le
demi-plan supérieur \(\mathbb H\) ; ce n’est ni la rotation \(J_2\), ni
l’échange \(S\), ni l’holonomie nonadique \(\Gamma_9\).

\begin{proposition}[Semi-conjugaison modulaire du rayon]
\label{prop:ii-t-modular}
Soient \(\iota:\mathbb R_{>0}\to\mathbb H\),
\(\iota(R_0)=iR_0^2\), et
\(S_{\rm mod}(\tau)=-1/\tau\). Alors
\begin{equation}
 S_{\rm mod}\circ\iota(R_0)=\iota(R_0^{-1}),\qquad
 J(iR_0^2)=J(iR_0^{-2}).
 \label{eq:ii-t-modular}
\end{equation}
En particulier, les deux formes étiquetées de l’ellipse ont la même
valeur de \(J\).
\end{proposition}

\begin{proof}
La première identité est
\(-1/(iR_0^2)=iR_0^{-2}\). La matrice
\(\left(\begin{smallmatrix}0&-1\\1&0\end{smallmatrix}\right)\)
appartient à \(\mathrm{SL}_2(\mathbb Z)\) et agit sur
\(\mathbb H\) comme \(S_{\rm mod}\). L’invariance modulaire de
\(J\), établie dans sa construction exceptionnelle, donne
\(J(S_{\rm mod}\tau)=J(\tau)\). En l’appliquant à
\(\tau=\iota(R_0)\), on obtient la seconde identité.
La proposition~\ref{prop:ii-t-radio} identifie ces deux arguments
aux modules des ellipses de formes opposées.
\end{proof}

La composition conserve ainsi une chaîne explicite :
\[
 (A,C^*,\hbar)\longmapsto(a_S,b_S)
 \longmapsto R_0\longmapsto iR_0^2\longmapsto J(iR_0^2).
\]
L’inversion de forme se transporte en inversion de rayon et en
\(S_{\rm mod}\) ; l’action orientée et le changement de section conservent
leurs applications propres. La fonction \(J\) assigne une valeur à une classe modulaire, tandis que
le marquage des axes distingue ses deux présentations et que l’état TPK
conserve sa généalogie. Telle est la relation démontrée entre l’ellipse
d’action, la dualité T et Moonshine : une composition avec des domaines et des
identités de transport, dont la valeur modulaire ne remplace ni l’opérateur
réticulaire ni la construction de \(V^\natural\).
